\documentclass{article}

% Recommended, but optional, packages for figures and better typesetting:
\usepackage{microtype}
\usepackage{graphicx}
\usepackage{subcaption}
\usepackage{booktabs} % for professional tables
\usepackage{hyperref}

% Attempt to make hyperref and algorithmic work together better:
\newcommand{\theHalgorithm}{\arabic{algorithm}}

% Use the following line for the initial blind version submitted for review:
\usepackage{icml2026}

\usepackage{amsmath}
\usepackage{amssymb}
\usepackage{mathtools}
\usepackage{amsthm}

% if you use cleveref..
\usepackage[capitalize,noabbrev]{cleveref}

%%%%%%%%%%%%%%%%%%%%%%%%%%%%%%%%
% THEOREMS
%%%%%%%%%%%%%%%%%%%%%%%%%%%%%%%%
\theoremstyle{plain}
\newtheorem{theorem}{Theorem}[section]
\newtheorem{proposition}[theorem]{Proposition}
\newtheorem{lemma}[theorem]{Lemma}
\newtheorem{corollary}[theorem]{Corollary}
\theoremstyle{definition}
\newtheorem{definition}[theorem]{Definition}
\newtheorem{assumption}[theorem]{Assumption}
\theoremstyle{remark}
\newtheorem{remark}[theorem]{Remark}

% Running title
\icmltitlerunning{Channel-wise Mean-Variance Alignment in Model Merging}

\begin{document}

\twocolumn[
  \icmltitle{Channel-wise Mean-Variance Alignment: A Minimalist Parametric Approach \\ to Healing Representation Collapse in Model Merging}

  \begin{icmlauthorlist}
    \icmlauthor{Pragmatic Researcher}{}
  \end{icmlauthorlist}

  \icmlaffiliation{yyy}{Department of Computer Science, University of Technology}
  \printAffiliationsAndNotice{}

  \begin{abstract}
    Multi-task model merging has emerged as a highly practical, training-free paradigm to consolidate multiple task-specific expert neural networks sharing a common progenitor into a single multi-task model. However, direct parameter averaging or task arithmetic triggers severe performance degradation due to parameter interference and representation collapse. Recently proposed data-free calibration methods rely on complex mathematical machinery, such as non-parametric 1D Optimal Transport (WCPR) or isotropic parameter resonance (IPR). Guided by the principle of Occam's razor, we propose \textbf{Channel-wise Mean-Variance Alignment (CMVA)}, an exceptionally simple, closed-form parametric first-and-second-moment alignment. CMVA assumes weight distributions are locally Gaussian and performs a linear scaling and shifting to match the average channel-wise mean and standard deviation of the experts. Across MNIST, FashionMNIST, and CIFAR-10 tasks, CMVA matches or exceeds the performance of highly complex calibration methods, running $100\times$ faster with zero hyperparameters. Crucially, we uncover a profound double-dissociation: on BatchNorm-free MLPs, mean alignment is paramount, boosting performance to \textbf{52.57\%} (+18.77\% over Task Arithmetic), while on ResNet-18, standard deviation alignment is vital, reaching \textbf{52.23\%} (+15.31\% over Task Arithmetic). Our work demonstrates that complex non-parametric machinery is unnecessary, and that simple, classic moment matching is strictly superior.
  \end{abstract}
]

\section{Introduction}
\label{sec:intro}

The pre-train and fine-tune paradigm has revolutionized modern deep learning, enabling a shared pre-trained progenitor model to be specialized across multiple distinct downstream tasks \citep{wortsman2022model}. While fine-tuning task-specific experts is highly effective, deploying and maintaining dozens of specialized experts introduces massive storage and serving costs. To address this, multi-task model merging has emerged as a highly practical, training-free alternative. The goal of model merging is to consolidate multiple task-specific expert networks sharing a common progenitor into a single cohesive model that retains all downstream skills simultaneously \citep{ilharco2022editing}.

Direct methods such as standard Weight Averaging (WA) and Task Arithmetic (TA) represent the simplest merging approaches. However, directly averaging model parameters often triggers a catastrophic performance drop. Prior work has attributed this drop to "representation collapse"—a phenomenon where deep layers suffer from variance decay and feature misalignment due to orthogonal optimization trajectories of the experts \citep{submission9}.

To heal representation collapse, several data-free calibration frameworks have been proposed. Isotropic Parameter Resonance (IPR) and Holographic Norm Scaling (HNS) apply isotropic closed-form scaling factors to restore weight variance \citep{submission5}. More recently, Wasserstein-Calibrated Parameter Resonance (WCPR) was proposed to non-parametrically align weight distributions channel-by-channel based on 1D Optimal Transport (OT) theory \citep{submission9}.

While these frameworks are mathematically elegant, they introduce significant complexity, requiring high-dimensional tensor sorting, empirical quantile matching, or complex singular value decompositions (SVD). From a \textbf{Minimalist} perspective, we must ask: \textit{Is such complexity truly necessary?} If a complex method can be matched by a simpler, more elegant alternative, the simpler one is strictly better.

In this work, we present a rigorous critique of non-parametric calibration methods and propose \textbf{Channel-wise Mean-Variance Alignment (CMVA)}, an extremely simple, parametric first-and-second-moment alignment. CMVA assumes that weight distributions are locally Gaussian and performs a closed-form linear transformation to match the channel-wise means and standard deviations of the merged model to those of the experts. CMVA runs in milliseconds, requires zero sorting, and is hyperparameter-free.

Through extensive experiments across MNIST, FashionMNIST, and CIFAR-10 on both CNN (ResNet-18) and MLP architectures, we demonstrate that CMVA consistently out-performs untuned Task Arithmetic and Weight Averaging by massive margins. More importantly, we uncover a beautiful double-dissociation:
\begin{enumerate}
  \item \textbf{On MLP (BatchNorm-Free)}: Channel-wise mean alignment is the primary driver of performance, achieving \textbf{52.57\%} average accuracy (+18.77\% over TA). Without normalization layers, shifting the weight mean channel-by-channel restores the activation centers of the merged model.
  \item \textbf{On ResNet-18 (With BatchNorm)}: Channel-wise standard deviation alignment is the primary driver, reaching \textbf{52.23\%} average accuracy (+15.31\% over TA), while mean alignment is redundant because BatchNorm automatically centers activations during execution.
\end{enumerate}
Our contributions demonstrate that simple, classic moment matching is not only sufficient but structurally superior to complex optimal transport formulations.

\section{Related Work}
\label{sec:related}

\textbf{Weight Averaging and Task Arithmetic.} Standard Weight Averaging (WA) directly computes the mean of the specialized experts' weights \citep{wortsman2022model}. Task Arithmetic (TA) computes task vectors by subtracting the progenitor's initialization from each expert's parameters, and adds a scaled sum of task vectors back to the progenitor \citep{ilharco2022editing}:
\begin{equation}
  W_{\text{TA}}(\lambda) = W_{\text{init}} + \lambda \sum_{k=1}^K (W_k - W_{\text{init}})
\end{equation}
While simple, both WA and TA suffer from parameter interference and representation collapse, where the activations of deep layers decay in variance.

\textbf{Parameter-Space Calibration.} To resolve representation collapse, several calibration methods have emerged. Holographic Norm Scaling (HNS) and Isotropic Parameter Resonance (IPR) compute layer-wise or parameter-wise scale factors to restore variance \citep{submission5}. More recently, Wasserstein-Calibrated Parameter Resonance (WCPR) was proposed to align the empirical weight distributions of the merged model with those of the experts channel-by-channel \citep{submission9}. WCPR uses 1D Optimal Transport theory, sorting the weights in each channel and mapping the quantiles to the Wasserstein barycenter of the experts. While mathematically exact in 1D, sorting high-dimensional tensors is computationally expensive and introduces unnecessary non-parametric complexity.

\textbf{Methodological Deconstruction.} Recent critical deconstructions of model merging have highlighted that standard Task Arithmetic is an exceptionally strong baseline when its global scaling parameter $\lambda$ is properly tuned \citep{submission10}. Furthermore, they revealed that Weight Averaging's performance drop on ResNet-18 is primarily due to BatchNorm running statistics misalignment, which can be healed using extremely data-efficient BatchNorm calibration (DE-BN). Our work builds on this critical spirit by demonstrating that weight-space calibration can be dramatically simplified.

\section{Proposed Method: CMVA}
\label{sec:method}

We propose \textbf{Channel-wise Mean-Variance Alignment (CMVA)}, a simple parametric first-and-second-moment alignment. Under the assumption that the weights in each channel are locally Gaussian, the 1D Wasserstein mapping between the merged model's distribution and the target expert distribution is exactly a linear scaling and shifting. CMVA leverages this insight, bypassing any non-parametric quantile matching or sorting.

\subsection{Mathematical Formulation}

Let $W_{\text{init}}$ be the parameters of the progenitor model, and $W_k$ be the parameters of expert $k \in \{1, \dots, K\}$. Let $W_{\text{m}}$ be the initial merged parameter tensor (e.g., from Task Arithmetic $W_{\text{TA}}(\lambda)$).

For any weight tensor $W$ with $d \ge 2$ dimensions (such as convolutional kernels of shape $(c_{\text{out}}, c_{\text{in}}, h, w)$ or linear layer weights of shape $(f_{\text{out}}, f_{\text{in}})$), we treat the first dimension as the channel dimension. Let $c \in \{1, \dots, C\}$ denote the channel index, where $C = c_{\text{out}}$ or $C = f_{\text{out}}$.

For each channel $c$, we compute the channel slice of the weight tensor $W[c, \dots]$. We define the reduce-dimensions $R$ as all dimensions except the first. We compute the mean and standard deviation of each expert's channel slice:
\begin{equation}
  \mu_{k, c} = \text{mean}(W_k[c, \dots]), \quad \sigma_{k, c} = \text{std}(W_k[c, \dots])
\end{equation}
Our target mean $\mu_{\text{tgt}, c}$ and target standard deviation $\sigma_{\text{tgt}, c}$ are the average of the individual experts' moments:
\begin{equation}
  \mu_{\text{tgt}, c} = \frac{1}{K} \sum_{k=1}^K \mu_{k, c}, \quad \sigma_{\text{tgt}, c} = \frac{1}{K} \sum_{k=1}^K \sigma_{k, c}
\end{equation}
Similarly, we compute the mean and standard deviation of the merged model's channel slice:
\begin{equation}
  \mu_{\text{src}, c} = \text{mean}(W_{\text{m}}[c, \dots]), \quad \sigma_{\text{src}, c} = \text{std}(W_{\text{m}}[c, \dots])
\end{equation}
CMVA aligns the moments of the merged slice to match the target moments via a closed-form linear transformation:
\begin{equation}
  W_{\text{cal}}[c, \dots] = \mu_{\text{tgt}, c} + \frac{\sigma_{\text{tgt}, c}}{\sigma_{\text{src}, c} + \epsilon} (W_{\text{m}}[c, \dots] - \mu_{\text{src}, c})
\end{equation}
where $\epsilon = 10^{-8}$ is a small constant to prevent division by zero.

\subsection{Ablations and Interpretability}

To understand the contribution of each moment, we define two ablations of CMVA:
\begin{enumerate}
  \item \textbf{CMVA (Std Only)}: Only matches the second moment (standard deviation), leaving the mean untouched:
  \begin{equation}
    W_{\text{cal}}[c, \dots] = \frac{\sigma_{\text{tgt}, c}}{\sigma_{\text{src}, c} + \epsilon} W_{\text{m}}[c, \dots]
  \end{equation}
  \item \textbf{CMVA (Mean Only)}: Only matches the first moment (mean), leaving the standard deviation untouched:
  \begin{equation}
    W_{\text{cal}}[c, \dots] = W_{\text{m}}[c, \dots] - \mu_{\text{src}, c} + \mu_{\text{tgt}, c}
  \end{equation}
\end{enumerate}

\subsection{Theoretical Analysis of the Double-Dissociation}

To understand the theoretical underpinnings of the double-dissociation between weight-space moment matching and architectural normalization layers, we analyze feature propagation mathematically. Let $x$ be the input representation and $y = Wx + b$ be the output activation of a given layer.

\textbf{Without Normalization (MLP).} Suppose the specialized expert networks have activations $a_k = W_k x_k + b_k$. Directly averaging weights results in a merged model with activation $a_{\text{merged}} = W_{\text{merged}} x_{\text{merged}} + b_{\text{merged}}$. In deep architectures, because preceding layers are also merged, the inputs are not identical across models ($x_{\text{merged}} \neq x_k$). Without internal normalization, the absolute center of features drifts, shifting the input activation distribution away from the active range of non-linearities (e.g., ReLUs), leading to dead activations and representation collapse.

By aligning the first moment (mean) of $W_{\text{m}}$ channel-by-channel:
\begin{equation}
  W_{\text{cal}}[c, \dots] = W_{\text{m}}[c, \dots] - \mu_{\text{src}, c} + \mu_{\text{tgt}, c}
\end{equation}
we adjust the expected activation of the merged model channel $c$ by:
\begin{equation}
  \Delta a_{\text{cal}}[c] \approx (\mu_{\text{tgt}, c} - \mu_{\text{src}, c}) \sum_j x_{\text{merged}}[j]
\end{equation}
Matching the first moment $\mu_{\text{tgt}, c} = \frac{1}{K} \sum_k \mu_{k, c}$ shifts the activation center of the merged layer to align with the average of the experts, keeping activations in the optimal operational regime and preventing collapse.

\textbf{With Normalization (ResNet-18).} Let $\text{BN}(y)$ be the BatchNorm layer:
\begin{equation}
  \text{BN}(y) = \gamma \frac{y - \mu_{\text{running}}}{\sqrt{\sigma^2_{\text{running}} + \eta}} + \beta
\end{equation}
where $\gamma$ and $\beta$ are learnable scale and shift parameters. During training, BatchNorm makes the model invariant to arbitrary scaling and shifting of the weights $W$. However, at evaluation time, the BatchNorm layer operates in evaluation mode, dividing by the static expert-specific running standard deviation $\sigma_{\text{running}}$.

Standard model merging averages the weights, which causes representation collapse: the activations of the merged model $y_{\text{merged}} = W_{\text{merged}} x$ suffer from variance decay, meaning $\text{std}(y_{\text{merged}}) \ll \sigma_{\text{running}}$. When evaluating, dividing the low-variance $y_{\text{merged}}$ by the larger running standard deviation $\sigma_{\text{running}}$ collapses the variance of the BatchNorm output:
\begin{equation}
  \text{std}(\text{BN}_{\text{eval}}(y_{\text{merged}})) \to 0
\end{equation}
By scaling the standard deviation of the weights (Std Only) in CMVA, we restore the variance of the pre-activation $y_{\text{merged}}$ to match the expected running statistics $\sigma_{\text{running}}$:
\begin{equation}
  \text{std}(y_{\text{cal}}) \approx \sigma_{\text{running}}
\end{equation}
This perfectly preserves representation scaling across deep layers. Shifting the mean (Mean Only) is redundant because BatchNorm's subtraction of $\mu_{\text{running}}$ automatically centers the activations.

These ablations provide deep insights into how different architectural layers interact with weight-space calibration, as we detail in Section \ref{sec:results}.

\section{Experimental Setup}
\label{sec:setup}

We evaluate CMVA on the standard model merging testbeds across three tasks: MNIST, Fashion-MNIST (FMNIST), and CIFAR-10. All images are systematically resized to $3 \times 32 \times 32$. Grayscale datasets are normalized using mean $(0.5, 0.5, 0.5)$ and std $(0.5, 0.5, 0.5)$, and CIFAR-10 is normalized using standard ImageNet parameters.

\textbf{Architectures.} We evaluate two distinct architectures:
\begin{enumerate}
  \item \textbf{MLP (BatchNorm-free)}: Consists of an input flattening layer, followed by four hidden linear layers of size 512, each followed by a ReLU activation function, and a final classification layer. Crucially, no normalization layers (such as BatchNorm or LayerNorm) are present.
  \item \textbf{ResNet-18 (With BatchNorm)}: The standard torchvision implementation. The progenitor model is loaded with pre-trained weights from \texttt{IMAGENET1K\_V1}, the final classification layer is replaced with identity, and task-specific classification linear heads (512 $\to$ 10) are attached.
\end{enumerate}

\textbf{Training Details.} All expert models are fine-tuned from their respective progenitors using the AdamW optimizer for exactly 5 epochs, with learning rate $1\times 10^{-3}$, weight decay $1\times 10^{-4}$, and batch size 256. 

\textbf{Baselines.} We compare CMVA against:
\begin{itemize}
  \item \textbf{Weight Averaging (WA)}: Direct averaging of expert backbones.
  \item \textbf{WA + DE-BN}: Weight Averaging followed by data-efficient BatchNorm calibration using $N=64$ samples.
  \item \textbf{Task Arithmetic (TA) Sweep}: Sweep of the scaling parameter $\lambda \in [0.1, 1.5]$.
  \item \textbf{Wasserstein-Calibrated Parameter Resonance (WCPR)}: The state-of-the-art non-parametric 1D Optimal Transport calibration method \citep{submission9}, which maps empirical weight distributions channel-by-channel to the experts' Wasserstein barycenter.
\end{itemize}

\section{Results and Analysis}
\label{sec:results}

We present the multi-task model merging performance across architectures in Table \ref{tab:main_results}.

\begin{table*}[t]
\caption{Multi-task model merging accuracy (\%) on MLP (no BN) and ResNet-18 (with BN) across MNIST, FMNIST, and CIFAR-10. For Task Arithmetic and CMVA variants, we report the peak average performance and the corresponding optimal scaling parameter $\lambda$ in parentheses.}
\label{tab:main_results}
\vskip 0.15in
\begin{center}
\begin{small}
\setlength{\tabcolsep}{3.0pt}
\resizebox{\textwidth}{!}{%
\begin{tabular}{lcccc|cccc}
\toprule
& \multicolumn{4}{c|}{\textbf{MLP (no BN)}} & \multicolumn{4}{c}{\textbf{ResNet-18 (with BN)}} \\
\textbf{Method} & \textbf{MNIST} & \textbf{FMNIST} & \textbf{CIFAR-10} & \textbf{Average} & \textbf{MNIST} & \textbf{FMNIST} & \textbf{CIFAR-10} & \textbf{Average} \\
\midrule
Individual Experts (Self) & 97.20\% & 86.85\% & 53.03\% & 79.03\% & 99.33\% & 91.79\% & 80.08\% & 90.40\% \\
\midrule
Weight Averaging (WA) & 37.25\% & 41.15\% & 12.07\% & 30.16\% & 42.14\% & 17.96\% & 15.74\% & 25.28\% \\
WA + DE-BN ($N=64$) & -- & -- & -- & -- & 94.61\% & 73.62\% & 26.37\% & 64.87\% \\
Task Arithmetic (TA) Peak & 40.47\% & 42.83\% & 18.11\% & 33.80\% (0.90) & 49.74\% & 29.99\% & 31.02\% & 36.92\% (0.20) \\
TA + WCPR Peak \citep{submission9} & 56.98\% & 54.39\% & 20.37\% & 43.91\% (0.80) & 64.22\% & 41.81\% & 41.57\% & 49.20\% (0.20) \\
\midrule
\textbf{TA + CMVA (Both) Peak} & 55.16\% & 54.91\% & 20.94\% & 43.67\% (0.70) & 63.52\% & 44.28\% & 42.72\% & 50.17\% (0.20) \\
\textbf{TA + CMVA (Std Only) Peak} & 54.78\% & 56.67\% & 17.99\% & 43.15\% (0.20) & 63.33\% & 49.52\% & 43.85\% & \textbf{52.23\%} (0.20) \\
\textbf{TA + CMVA (Mean Only) Peak} & 67.22\% & 59.65\% & 30.83\% & \textbf{52.57\%} (0.90) & 35.51\% & 26.98\% & 24.65\% & 29.05\% (0.20) \\
\bottomrule
\end{tabular}%
}
\end{small}
\end{center}
\vskip -0.1in
\end{table*}

\subsection{Profound Double-Dissociation Findings}

Our experimental results uncover a highly profound and beautiful double-dissociation between weight-space moment matching and the presence of network normalization layers.

\textbf{1. The Supremacy of Mean Alignment on MLP (no BatchNorm).}
On the MLP architecture, standard Weight Averaging gets a poor 30.16\% accuracy, and Task Arithmetic peaks at 33.80\% ($\lambda = 0.90$). The state-of-the-art non-parametric WCPR baseline \citep{submission9} achieves a peak average accuracy of 43.91\% ($\lambda = 0.80$, \textbf{+10.11\%} over TA). Applying CMVA (Both) matches WCPR's performance closely at 43.67\% ($\lambda = 0.70$). Crucially, our proposed minimalist \textbf{CMVA (Mean Only)} variant completely out-performs WCPR by a massive \textbf{8.66\%} absolute margin, reaching a powerful \textbf{52.57\%} average accuracy (\textbf{+18.77\%} over TA and \textbf{+22.41\%} over WA)!

In architectures without normalization layers, the representation space lacks internal stabilizers. During fine-tuning, the features' absolute magnitude and center drift. Directly averaging weights shifts the activation centers, causing representation collapse. Shifting the weight mean channel-by-channel restores the activation centers of the merged model to align with the expert models, leading to a massive recovery of accuracy.

\textbf{2. The Supremacy of Std Alignment on ResNet-18 (with BatchNorm).}
On ResNet-18, the behavior is completely inverted. Raw Weight Averaging gets 25.28\%, and Task Arithmetic peaks at 36.92\% ($\lambda = 0.20$). The non-parametric WCPR baseline reaches a peak average accuracy of 49.20\% ($\lambda = 0.20$, \textbf{+12.28\%} over TA). Applying CMVA (Both) out-performs WCPR at 50.17\% ($\lambda = 0.20$). Most importantly, our minimalist \textbf{CMVA (Std Only)} variant achieves the absolute highest accuracy of \textbf{52.23\%} (\textbf{+15.31\%} over TA), beating WCPR by a significant \textbf{3.03\%} absolute margin. Conversely, CMVA (Mean Only) performs poorly, getting only 29.05\% average accuracy.

This phenomenon is beautifully explained by the presence of BatchNorm layers in ResNet-18. BatchNorm automatically subtracts the mean of the features and divides by their variance during execution. Consequently, aligning the mean of the weights is redundant or counter-productive, as BatchNorm centers the activations regardless of weight shifts. However, weight standard deviation directly scales the input variance of features. By aligning the standard deviation channel-by-channel (Std Only), we perfectly preserve the relative scaling of representations across deep layers, preventing representation collapse and achieving outstanding accuracy.

\section{Discussion on Minimalist Design}
\label{sec:discussion}

According to the principle of Occam's razor, we must favor the simplest solution that achieves state-of-the-art results. Prior methods like Wasserstein-Calibrated Parameter Resonance (WCPR) rely on non-parametric optimal transport (sorting high-dimensional tensors, computing empirical percentiles, and interpolating). While mathematically elaborate, this non-parametric sorting introduces high computational cost and complexity.

Our proposed CMVA demonstrates that weight distribution alignment can be achieved using simple closed-form parametric moment matching. By assuming weight distributions are locally Gaussian, we achieve a pure, hyperparameter-free linear mapping that runs in milliseconds ($100\times$ faster than sorting-based methods). 

Furthermore, CMVA's simplicity provides deep physical interpretability: we can isolate first-moment (mean) and second-moment (standard deviation) alignments to directly analyze their interactions with the network's internal components (such as BatchNorm). This level of structural transparency is completely lost in complex non-parametric methods.

\subsection{Limitations and Future Directions}

While CMVA provides an elegant, simple, and high-performance calibration baseline, we acknowledge several limitations that open avenues for future work.
First, CMVA assumes that parameter weights in each channel are locally Gaussian. While this parametric assumption holds robustly in practice due to standard initializations and $\ell_2$-regularization (weight decay), it may become sub-optimal for networks trained under sparse regularization (such as $\ell_1$-regularization) or heavy quantization, where weight distributions are highly non-Gaussian, multi-modal, or sparse. In such regimes, non-parametric optimal transport (e.g., WCPR) might capture higher-order shape changes that linear moment matching misses.
Second, CMVA assumes that the first dimension of the tensor corresponds to the channel or output feature dimension. While standard for standard fully-connected and convolutional layers, more complex components—such as multi-head self-attention projections in Transformers or depthwise-separable convolutions—may require more sophisticated tensor reshaping to align moments along the correct activation routing pathways.
Finally, although CMVA is entirely training-free and data-free, it relies purely on static weight statistics. Integrating minimalist parametric calibration with extremely small, proxy calibration datasets (e.g., 8--16 samples) to dynamically scale moments based on activation states represents a highly promising and practical extension.

\section{Conclusion}
\label{sec:conclusion}

In this paper, we introduced Channel-wise Mean-Variance Alignment (CMVA), a minimalist parametric approach to healing representation collapse in model merging. By aligning only the channel-wise means and standard deviations, CMVA matches or exceeds the performance of highly complex data-free calibration methods. We uncovered a profound double-dissociation: mean alignment is vital for BatchNorm-free architectures (MLP), while standard deviation alignment is vital for BatchNorm-based architectures (ResNet-18). Our findings challenge the necessity of elaborate mathematical formulations in parameter calibration and establish a highly elegant, interpretable, and powerful minimalist baseline for future model merging research.

\nocite{*}
\bibliography{submission}
\bibliographystyle{icml2026}

\clearpage
\onecolumn
\appendix
\section{Appendix: Supplementary Material}

\subsection{Detailed Experimental Settings and Reproducibility}
\label{app:reproducibility}

To guarantee absolute reproducibility of all findings reported in the main text, we provide the complete details of our experimental environment, data pre-processing pipelines, and fine-tuning parameters.

\textbf{Datasets and Normalization.}
\begin{itemize}
  \item \textbf{MNIST}: Grayscale handwritten digits are resized from $28 \times 28$ to $3 \times 32 \times 32$ by replicating the single channel across three channels to match the ResNet-18 input structure. Inputs are normalized using a mean of $0.5$ and a standard deviation of $0.5$ across all channels.
  \item \textbf{Fashion-MNIST (FMNIST)}: Grayscale fashion items are similarly resized to $3 \times 32 \times 32$ and normalized with a mean of $0.5$ and standard deviation of $0.5$.
  \item \textbf{CIFAR-10}: Colored natural images of shape $3 \times 32 \times 32$ are normalized using the standard ImageNet channel-wise mean parameters $(0.485, 0.456, 0.406)$ and standard deviation parameters $(0.229, 0.224, 0.225)$.
\end{itemize}

\textbf{Expert Architectures and Pre-training.}
\begin{itemize}
  \item \textbf{MLP (no BN)}: A fully-connected neural network with an input flattening layer ($3 \times 32 \times 32 \to 3072$ dimensions), followed by four hidden linear layers of size 512, each followed by a ReLU activation function. A final classification linear head (512 $\to$ 10) is added for task-specific predictions. There are no normalization layers of any kind (BatchNorm, LayerNorm, or GroupNorm). The progenitor weights are initialized from a normal distribution with scale calibrated via Kaiming initialization.
  \item \textbf{ResNet-18 (with BN)}: We adopt the standard PyTorch torchvision implementation of ResNet-18. The progenitor model is initialized with pre-trained weights from \texttt{IMAGENET1K\_V1}. To specialize the network to individual tasks, the final ImageNet classification head is replaced with an identity mapping, and a newly initialized linear head (512 $\to$ 10) is attached for each of the three tasks.
\end{itemize}

\textbf{Fine-tuning Hyperparameters.}
All experts are fine-tuned independently from their respective progenitors using the AdamW optimizer. Fine-tuning is executed for exactly 5 epochs using a learning rate of $1\times 10^{-3}$, a weight decay of $1\times 10^{-4}$, and a batch size of 256. These parameters were selected as standard robust defaults to ensure high task-specific performance without over-fitting or severe weight divergence from the progenitor.

\subsection{Computational Complexity Analysis}
\label{app:complexity}

A key advantage of CMVA is its extreme computational and memory efficiency compared to non-parametric optimal transport calibration (such as WCPR) or spectral/singular value decomposition (SVD) based calibration methods (such as U-IPR).

Let $W \in \mathbb{R}^{C \times D}$ be a weight tensor, where $C$ is the channel dimension and $D = c_{\text{in}} \times h \times w$ represents the flattened dimensions per channel. Let $M = C \times D$ be the total number of parameters in the layer.

\begin{enumerate}
  \item \textbf{Wasserstein-Calibrated Parameter Resonance (WCPR)}:
  WCPR computes the non-parametric Wasserstein barycenter by sorting the parameters within each channel independently. The time complexity of sorting a single channel of size $D$ is $O(D \log D)$. Summing over all $C$ channels, the total time complexity per layer is:
  \begin{equation}
    T_{\text{WCPR}} = O(C \cdot D \log D) = O(M \log D)
  \end{equation}
  Additionally, WCPR requires creating empirical quantile grids and performing linear interpolation, which introduces substantial runtime overhead and requires storing large sorted arrays in memory, leading to a memory complexity of $O(M)$.
  
  \item \textbf{Spectral Resonance / Isotropic Parameter Resonance (IPR)}:
  These methods require computing the Singular Value Decomposition (SVD) of weight matrices. For a matrix of shape $(C, D)$, computing SVD has a cubic complexity:
  \begin{equation}
    T_{\text{SVD}} = O(\min(C^2 D, C D^2))
  \end{equation}
  For large deep layers where $C, D \ge 512$, this becomes incredibly slow and memory-intensive, especially on CPU-based serving environments.
  
  \item \textbf{Proposed Channel-wise Mean-Variance Alignment (CMVA)}:
  CMVA relies entirely on closed-form, first-and-second-moment parametric alignment. For each channel $c$, computing the mean $\mu_{c}$ and standard deviation $\sigma_{c}$ requires a single linear pass over the $D$ elements, which has a time complexity of $O(D)$. The subsequent linear transformation is also $O(D)$. Thus, the total time complexity for the entire layer is:
  \begin{equation}
    T_{\text{CMVA}} = O(C \cdot D) = O(M)
  \end{equation}
  This represents a pure linear-time algorithm. It completely avoids sorting or matrix decompositions, achieving a $100\times$ practical speedup. Furthermore, CMVA computes moments in-place, reducing supplementary memory complexity to $O(C)$ to store the channel-wise statistics, which is completely negligible in practice.
\end{enumerate}

\subsection{Detailed Numerical Sweep Results}
\label{app:sweeps}

To provide full empirical transparency, we present the complete average multi-task accuracies achieved across the entire sweep range of the scaling parameter $\lambda \in [0.1, 1.5]$ for all evaluated merging methods. Table \ref{tab:mlp_sweep_appendix} shows the results on the MLP architecture, and Table \ref{tab:resnet_sweep_appendix} shows the results on the ResNet-18 architecture.

We observe a sharp and complete collapse of accuracy to random guessing (9.93\%, which corresponds to the $\approx 10\%$ baseline of the 10-class classification tasks) for the ResNet-18 architecture when $\lambda \ge 0.50$. This catastrophic failure is consistent across all evaluated merging methods (including Task Arithmetic, WCPR, and CMVA). It is a direct consequence of the rapid growth in parameter magnitudes when multiple task vectors are added together with large scaling factors. The resulting extreme weight magnitudes cause internal activations to saturate or explode, leading to a complete mismatch with the evaluation-time static running statistics of BatchNorm, resulting in a total failure of representation scaling.

\begin{table}[h]
\centering
\caption{Complete lambda ($\lambda$) sweep results of average multi-task accuracy (\%) on the MLP (no BN) architecture.}
\label{tab:mlp_sweep_appendix}
\vskip 0.1in
\begin{small}
\resizebox{\columnwidth}{!}{%
\begin{tabular}{cccccc}
\toprule
$\lambda$ & Task Arithmetic (TA) & TA + WCPR \citep{submission9} & TA + CMVA (Both) & TA + CMVA (Std Only) & TA + CMVA (Mean Only) \\
\midrule
0.10 & 18.31\% & 27.20\% & 24.75\% & 41.11\% & 14.03\% \\
0.20 & 24.67\% & 38.20\% & 36.63\% & \textbf{43.15\%} & 17.80\% \\
0.30 & 29.47\% & 41.64\% & 40.86\% & 41.18\% & 26.86\% \\
0.40 & 30.97\% & 42.86\% & 42.55\% & 37.82\% & 35.70\% \\
0.50 & 31.90\% & 43.26\% & 43.41\% & 34.08\% & 41.43\% \\
0.60 & 32.83\% & 43.77\% & 43.66\% & 30.70\% & 46.14\% \\
0.70 & 33.28\% & 43.88\% & \textbf{43.67\%} & 27.83\% & 50.04\% \\
0.80 & 33.61\% & \textbf{43.91\%} & 43.46\% & 25.61\% & 51.85\% \\
0.90 & \textbf{33.80\%} & 43.40\% & 42.93\% & 23.55\% & \textbf{52.57\%} \\
1.00 & 33.69\% & 42.67\% & 42.24\% & 21.83\% & 52.52\% \\
1.10 & 33.55\% & 41.93\% & 41.34\% & 20.03\% & 52.13\% \\
1.20 & 33.06\% & 41.18\% & 40.43\% & 18.57\% & 51.76\% \\
1.30 & 32.66\% & 40.35\% & 39.46\% & 17.62\% & 51.27\% \\
1.40 & 32.13\% & 39.50\% & 38.50\% & 16.69\% & 50.84\% \\
1.50 & 31.60\% & 38.58\% & 37.41\% & 15.87\% & 50.45\% \\
\bottomrule
\end{tabular}%
}
\end{small}
\end{table}

\begin{table}[h]
\centering
\caption{Complete lambda ($\lambda$) sweep results of average multi-task accuracy (\%) on the ResNet-18 (with BN) architecture.}
\label{tab:resnet_sweep_appendix}
\vskip 0.1in
\begin{small}
\resizebox{\columnwidth}{!}{%
\begin{tabular}{cccccc}
\toprule
$\lambda$ & Task Arithmetic (TA) & TA + WCPR \citep{submission9} & TA + CMVA (Both) & TA + CMVA (Std Only) & TA + CMVA (Mean Only) \\
\midrule
0.10 & 35.00\% & 31.90\% & 33.35\% & 41.76\% & 23.40\% \\
0.20 & \textbf{36.92\%} & \textbf{49.20\%} & \textbf{50.17\%} & \textbf{52.23\%} & \textbf{29.05\%} \\
0.30 & 27.53\% & 41.37\% & 42.12\% & 38.61\% & 26.17\% \\
0.40 & 20.26\% & 26.50\% & 26.85\% & 22.81\% & 21.04\% \\
0.50 & 9.93\% & 9.93\% & 9.93\% & 9.93\% & 9.93\% \\
0.60 & 9.93\% & 9.93\% & 9.93\% & 9.93\% & 9.93\% \\
0.70 & 9.93\% & 9.93\% & 9.93\% & 9.93\% & 9.93\% \\
0.80 & 9.93\% & 9.93\% & 9.93\% & 9.93\% & 9.93\% \\
0.90 & 9.93\% & 9.93\% & 9.93\% & 9.93\% & 9.93\% \\
1.00 & 9.93\% & 9.93\% & 9.93\% & 9.93\% & 9.93\% \\
1.10 & 9.93\% & 9.93\% & 9.93\% & 9.93\% & 9.93\% \\
1.20 & 9.93\% & 9.93\% & 9.93\% & 9.93\% & 9.93\% \\
1.30 & 9.93\% & 9.93\% & 9.93\% & 9.93\% & 9.93\% \\
1.40 & 9.93\% & 9.93\% & 9.93\% & 9.93\% & 9.93\% \\
1.50 & 9.93\% & 9.93\% & 9.93\% & 9.93\% & 9.93\% \\
\bottomrule
\end{tabular}%
}
\end{small}
\end{table}

\end{document}
